% =============================================================================
%  Appendix -- back-matter shown AFTER the Thank-you slide.
%  Contains Acronyms and Glossary for committee reference.
% =============================================================================

\appendix

\section{Appendix}

\sectiondivider{Appendix}{Acronyms and glossary for reference}{sec:appendix}

% --- Subsection: Acronyms --------------------------------------------------
\subsection{Acronyms}

% ---- Acronyms page 1: Models, layers, training -----------------------------
\begin{frame}{Acronyms -- models, layers \& training}
\hypertarget{sub:appendix-acronyms}{}%
\setlength{\extrarowheight}{1pt}
\setlength{\tabcolsep}{4pt}
\footnotesize
\begin{columns}[T,onlytextwidth]
\begin{column}{0.48\textwidth}
\textbf{\color{tecdark}Recurrent \& sequence models}\\[2pt]
\begin{tabular}{@{}l@{\hspace{8pt}}p{0.66\linewidth}@{}}
\textbf{LSTM}    & Long Short-Term Memory \\
\textbf{BiLSTM}  & Bidirectional LSTM \\
\textbf{GRU}     & Gated Recurrent Unit \\
\textbf{BiGRU}   & Bidirectional GRU \\
\textbf{RNN}     & Recurrent Neural Network \\
\textbf{TCN}     & Temporal Convolutional Network \\
\textbf{CNN}     & Convolutional Neural Network \\
\textbf{Conv1D}  & one-dimensional convolutional layer \\
\textbf{ST-GCN}  & Spatio-Temporal Graph Convolutional Network \\
\textbf{GAP}     & Global Average Pooling \\
\textbf{PAF}     & Part Affinity Field (in OpenPose) \\
\textbf{BPTT}    & Backpropagation Through Time \\
\end{tabular}
\end{column}
\hspace{0.02\textwidth}
\begin{column}{0.48\textwidth}
\textbf{\color{tecdark}Optimizers, losses \& regularization}\\[2pt]
\begin{tabular}{@{}l@{\hspace{8pt}}p{0.66\linewidth}@{}}
\textbf{Adam}    & Adaptive Moment estimation (optimizer) \\
\textbf{SGD}     & Stochastic Gradient Descent (optimizer) \\
\textbf{XEnt}    & Cross-entropy (loss function) \\
\textbf{MSE}     & Mean Squared Error (loss function) \\
\textbf{L2}      & L2 weight regularization (ridge / weight decay) \\
\end{tabular}\\[10pt]
\textbf{\color{tecdark}Quality classes (1--10 $\to$ 5)}\\[2pt]
\begin{tabular}{@{}l@{\hspace{8pt}}l@{}}
\textbf{1--2}    & Very Low \\
\textbf{3--4}    & Low \\
\textbf{5--6}    & Medium \\
\textbf{7--8}    & High \\
\textbf{9--10}   & Excellent \\
\end{tabular}
\end{column}
\end{columns}
\end{frame}

% ---- Acronyms page 2: Evaluation, statistics, tooling ----------------------
\begin{frame}{Acronyms -- evaluation, statistics \& tooling}
\setlength{\extrarowheight}{1pt}
\setlength{\tabcolsep}{4pt}
\footnotesize
\begin{columns}[T,onlytextwidth]
\begin{column}{0.48\textwidth}
\textbf{\color{tecdark}Evaluation \& statistics}\\[2pt]
\begin{tabular}{@{}l@{\hspace{8pt}}p{0.66\linewidth}@{}}
\textbf{TP}      & True Positive \\
\textbf{FP}      & False Positive \\
\textbf{FN}      & False Negative \\
\textbf{TN}      & True Negative \\
\textbf{F1}      & F1-score (harmonic mean of precision and recall) \\
\textbf{F1$_m$}  & Macro F1 (unweighted mean of per-class F1) \\
\textbf{$\rho$}  & Spearman's rank correlation coefficient \\
\textbf{$\kappa$}& Cohen's kappa (inter-rater agreement) \\
\textbf{ICC}     & Intraclass Correlation Coefficient \\
\textbf{pp}      & percentage points (absolute, not relative) \\
\textbf{LOPO}    & Leave-One-Player-Out (cross-validation) \\
\textbf{ROC}     & Receiver Operating Characteristic \\
\textbf{DET}     & Detection Error Tradeoff \\
\end{tabular}
\end{column}
\hspace{0.02\textwidth}
\begin{column}{0.48\textwidth}
\textbf{\color{tecdark}Hardware, data \& tooling}\\[2pt]
\begin{tabular}{@{}l@{\hspace{8pt}}p{0.66\linewidth}@{}}
\textbf{FPS}     & Frames Per Second \\
\textbf{IMU}     & Inertial Measurement Unit \\
\textbf{JSON}    & JavaScript Object Notation \\
\textbf{CLI}     & Command-Line Interface \\
\textbf{GUI}     & Graphical User Interface \\
\textbf{GPU}     & Graphics Processing Unit \\
\textbf{API}     & Application Programming Interface \\
\end{tabular}\\[10pt]
\textbf{\color{tecdark}Reading the metrics}\\[2pt]
{\scriptsize
A model with accuracy $a$ and Spearman $\rho$ near zero can rank classes randomly; one with modest $a$ but high $\rho$ preserves the quality ordering even when the exact class is wrong -- the regime we target on an ordinal problem.}
\end{column}
\end{columns}
\end{frame}

% --- Subsection: Glossary --------------------------------------------------
\subsection{Glossary}

% ---- Glossary page 1: Sport & vision (6 entries) ---------------------------
\begin{frame}{Glossary -- sport \& vision}
\hypertarget{sub:appendix-glossary}{}%
\footnotesize
\begin{description}\setlength{\itemsep}{8pt}\setlength{\parsep}{0pt}
\item[\textbf{Amateur}]\hfill\\
A sport player who participates primarily for enjoyment, personal satisfaction, and skill development rather than for financial gain.

\item[\textbf{Bandeja}]\hfill\\
Paddle tennis stroke typically executed in response to a mid-height lob from the opponent. The ideal hitting zone is mid-court, at a height between a volley and a smash.

\item[\textbf{Paddle tennis} (padel)]\hfill\\
Racket sport that merges aspects of tennis and squash. Played primarily in doubles on an enclosed court roughly one-third the size of a tennis court, with playable walls integral to the game.

\item[\textbf{Biomechanics}]\hfill\\
Study of the mechanical laws governing the movement and structure of living organisms; here, the analysis of joint trajectories, forces, and timing that characterize human motion in sports.

\item[\textbf{Pose estimation}]\hfill\\
Computer vision technique that detects and tracks human body key-points (e.g., joints) from images or videos to estimate spatial positions and orientations.

\item[\textbf{OpenPose}]\hfill\\
Open-source real-time multi-person 2D pose-estimation framework that predicts body, hand, facial, and foot key-points using confidence maps and Part Affinity Fields.
\end{description}
\end{frame}

% ---- Glossary page 2: Data & training-time concepts (5 entries) ------------
\begin{frame}{Glossary -- data \& training-time concepts}
\footnotesize
\begin{description}\setlength{\itemsep}{8pt}\setlength{\parsep}{0pt}
\item[\textbf{Stratified split}]\hfill\\
Partition of a dataset into training and test subsets that preserves the proportion of samples per class within each subset -- used here to keep the five quality classes balanced across both splits.

\item[\textbf{Data augmentation}]\hfill\\
Synthesizing additional training samples from existing ones by applying small perturbations (e.g., Gaussian noise on coordinate trajectories) to enlarge the effective training set and improve generalization on small datasets.

\item[\textbf{Cross-entropy loss}]\hfill\\
Standard classification loss that compares the predicted probability distribution against the one-hot encoded true label. Penalizes confident wrong predictions more strongly than uncertain ones.

\item[\textbf{Dropout}]\hfill\\
Regularization that randomly disables a fraction of neurons during each training step, forcing the network not to rely on any single feature path. This thesis uses 0.40 and 0.35 in the two BiLSTM layers.

\item[\textbf{L2 regularization} (weight decay)]\hfill\\
Penalty added to the loss proportional to the squared magnitudes of the model weights. Pushes weights toward zero, smooths decision boundaries, and helps when training data is scarce.
\end{description}
\end{frame}

% ---- Glossary page 3: Evaluation & generalization concepts (6 entries) -----
\begin{frame}{Glossary -- evaluation \& generalization}
\footnotesize
\begin{description}\setlength{\itemsep}{7pt}\setlength{\parsep}{0pt}
\item[\textbf{Ordinal classification}]\hfill\\
Classification in which target classes follow a natural ordering (\textit{Very Low} $<$ \textit{Low} $<$ \textit{Medium} $<$ \textit{High} $<$ \textit{Excellent}). Adjacent-class errors are less wrong than distant-class errors, motivating rank-based metrics such as Spearman's $\rho$.

\item[\textbf{Confusion matrix}]\hfill\\
$K\!\times\!K$ matrix where entry $(i, j)$ records true-class-$i$ samples predicted as class $j$. Diagonal entries are correct; off-diagonal entries reveal where errors land.

\item[\textbf{Overfitting}]\hfill\\
A model fitting the training data very closely -- including noise and player-specific idiosyncrasies -- but failing to generalize. Visible here as a $\sim$70 pp gap between train and test accuracy.

\item[\textbf{Inter-rater reliability}]\hfill\\
Statistical agreement between two or more annotators evaluating the same content (e.g., Cohen's $\kappa$, ICC). Not computed in this thesis -- each clip was graded by one coach.

\item[\textbf{Leave-one-player-out} (LOPO)]\hfill\\
Cross-validation protocol in which all clips of one player are held out for testing per fold; gives a conservative estimate of generalization to truly unseen individuals.

\item[\textbf{Percentage points} (pp)]\hfill\\
Absolute difference between two percentages: $26.8\%\!\to\!36.1\%$ is a gain of $9.3$~pp, not a $35\%$ relative increase.
\end{description}
\end{frame}
